\begin{abstract}

A free-flow road network can devolve into a traffic jam with a local disturbance that pushes the density past a critical point, a phenomenon known as a sudden traffic jam. This instability in traffic flow has been studied in the past with different infrastucutre and vehicle based control apprpoaches. One notable approach called self-regulating cars learned a controller that prevents this by setting a desired speed per road segment using a central cloud observer. In this paper, we ask what is the minimal agumentation needed to an ordinary vehicle for it to become self-regulating in a decentralized manner.

Our answer has three parts: a phone the driver already owns (OnePlus Nord N10), an onboard computer (NVIDIA Jetson Orin Nano- US\$250), and a public traffic API (HERE free tier). Nothing is installed on the road, and no server sits inside the control loop. The policy is a deterministic function of a feed every vehicle can query, so vehicles carrying the same weights compute the same action without coordinating. The phone samples its camera, GPS and inertial unit at rates the Jetson sets at run time rather than at install time.

Eight drives on one day recorded 22{,}929 ticks over 152.8 minutes on vehicle power. The on-Jetson path has a median of 19.8\,ms and the phone-to-Jetson path a 95th percentile of 116.19\,ms against a 200\,ms target. Nothing was connected to the vehicle's controls, and no advisory was acted on by a driver. In simulation on a prior-tested network, the deployed configuration gains $230 \pm 44$\,veh/h over no control while the published self-regulating cars configuration gains $224 \pm 61$\,veh/h, establishing that our decentralized deployment is as performant as the theoretical centralized controller. Two deployment failures are reported because a simulator cannot produce them: the phone overheated on a windscreen mount in sunlight and the android platform stalled the pipeline, and a ninety-degree camera rotation sent detection to a null during the drive.
\end{abstract}
